\section{Prayers of the Season}\label{sec:prayers}

Some of these prayers belong to every day of Advent: the Angelus, the
canticles of Zachary and of Mary, and the antiphon of Our Lady. Others belong to
particular days, from the responsory of the First Sunday to the great antiphons
of 17 to 23 December, and the last are for Communion, for confession and for
Christmas.

\subsection{The Angelus}

Said morning, noon and evening. Its prayer, \latin{Gratiam tuam}, is the
prayer of the Fourth Sunday of Advent (\ref{sec:fourth-sunday}).

\noindent \latin{V. Angelus Domini nuntiavit Mariae.}\par
\noindent \latin{R. Et concepit de Spiritu Sancto. Ave Maria.}\par
\noindent \latin{V. Ecce ancilla Domini.}\par
\noindent \latin{R. Fiat mihi secundum verbum tuum. Ave Maria.}\par
\noindent \latin{V. Et Verbum caro factum est.}\par
\noindent \latin{R. Et habitavit in nobis. Ave Maria.}\par
\noindent \latin{V. Ora pro nobis, sancta Dei Genitrix.}\par
\noindent \latin{R. Ut digni efficiamur promissionibus Christi.}\par
\noindent \latin{Oremus.}\par
\noindent \latin{Gratiam tuam, quaesumus, Domine, mentibus nostris infunde: ut qui, angelo nuntiante, Christi Filii tui Incarnationem cognovimus, per Passionem ejus et Crucem ad resurrectionis gloriam perducamur. Per eundem Christum Dominum nostrum. R. Amen.}\par

\medskip

\noindent V. The angel of the Lord declared unto Mary.\par
\noindent R. And she conceived of the Holy Ghost. Ave Maria.\par
\noindent V. Behold the handmaid of the Lord.\par
\noindent R. Be it done unto me according to thy word. Ave Maria.\par
\noindent V. And the Word was made Flesh.\par
\noindent R. And dwelt among us. Ave Maria.\par
\noindent V. Pray for us, O Holy Mother of God.\par
\noindent R. That we may be made worthy of the promises of Christ.\par
\noindent Let us pray.\par
\noindent Pour forth, we beseech Thee, O Lord, thy grace into our hearts; that we, to whom the Incarnation of Christ thy Son was made known by the message of an angel, may by his Passion and Cross be brought to the glory of his resurrection. Through the same Christ our Lord. R. Amen.\par

\subsection{The Great O Antiphons}

From 17 to 23 December the Church sings one of these antiphons each evening
at Evening Prayer, with the Magnificat, in the order of the dates below. At
Mass the verse before the Gospel takes them up, shortened and in another order
(\ref{sec:last-days}).

\subsubsection*{17 December}

\begin{quote}
\latin{O Sapientia, quae ex ore Altissimi prodiisti, attingens a fine usque ad finem, fortiter suaviterque disponens omnia: veni ad docendum nos viam prudentiae.}\par
\end{quote}

\begin{quote}
O Wisdom, That comest out of the mouth of the Most High, That reachest from one end to another, and dost mightily and sweetly order all things: come, to teach us the way of prudence!\par
\end{quote}

\subsubsection*{18 December}

\begin{quote}
\latin{O Adonai, et Dux domus Israël, qui Moysi in igne flammae rubi apparuisti, et ei in Sina legem dedisti: veni ad redimendum nos in brachio extento.}\par
\end{quote}

\begin{quote}
O Adonai, and Ruler of the house of Israël, Who didst appear unto Moses in the burning bush, and gavest him the law in Sinai: come, to redeem us with an outstretched arm!\par
\end{quote}

\subsubsection*{19 December}

\begin{quote}
\latin{O radix Jesse, qui stas in signum populorum, super quem continebunt reges os suum, quem Gentes deprecabuntur: veni ad liberandum nos, jam noli tardare.}\par
\end{quote}

\begin{quote}
O Root of Jesse, Which standest for an ensign of the people, at Whom the kings shall shut their mouths, unto Whom the Gentiles shall seek: come, to deliver us, make no tarrying!\par
\end{quote}

\subsubsection*{20 December}

\begin{quote}
\latin{O clavis David, et sceptrum domus Israël; qui aperis, et nemo claudit; claudis, et nemo aperit: veni, et educ vinctum de domo carceris, sedentem in tenebris, et umbra mortis.}\par
\end{quote}

\begin{quote}
O Key of David, and Sceptre of the house of Israël; That openest, and no man shutteth; and shuttest and no man openeth: come, to bring out the prisoners from the prison, and them that sit in darkness, and in the shadow of death!\par
\end{quote}

\subsubsection*{21 December}

\begin{quote}
\latin{O Oriens, splendor lucis aeternae, et sol justitiae: veni, et illumina sedentes in tenebris, et umbra mortis.}\par
\end{quote}

\begin{quote}
O Day-Spring, Brightness of the everlasting Light, Sun of Righteousness, come, to give light to them that sit in darkness, and in the shadow of death!\par
\end{quote}

\subsubsection*{22 December}

\begin{quote}
\latin{O Rex Gentium, et desideratus earum, lapisque angularis, qui facis utraque unum: veni, et salva hominem, quem de limo formasti.}\par
\end{quote}

\begin{quote}
O King of the Gentiles, yea, and Desire thereof, O Corner-stone that makest of twain one: come, to save man, whom Thou hast made of the dust of the earth!\par
\end{quote}

\subsubsection*{23 December}

\begin{quote}
\latin{O Emmanuel, Rex et legifer noster, exspectatio Gentium, et Salvator earum: veni ad salvandum nos Domine Deus noster.}\par
\end{quote}

\begin{quote}
O Emmanuel, our King and our Law-giver, Longing of the Gentiles, yea, and Salvation thereof: come to save us, O Lord our God!\par
\end{quote}

\subsection{\latin{Aspiciens a longe}}

The first responsory of Matins on the First Sunday of Advent, as the
Roman Breviary gave it. The present Office of Readings of that Sunday, after
St Cyril of Jerusalem's teaching on the two comings of Christ, sings a
responsory built on it.

\noindent \latin{R. Aspiciens a longe, ecce video Dei potentiam venientem, et nebulam totam terram tegentem. * Ite obviam ei, et dicite: * Nuntia nobis si tu es ipse, * Qui regnaturus es in populo Israël.}\par
\noindent \latin{V. Quique terrigenae, et filii hominum, simul in unum dives et pauper. Ite obviam ei, et dicite.}\par
\noindent \latin{V. Qui regis Israël, intende, qui deducis velut ovem Joseph. Nuntia nobis, si tu es ipse.}\par
\noindent \latin{V. Tollite portas principes vestras, et elevamini portae aeternales, et introibit Rex gloriae. Qui regnaturus es in populo Israël.}\par
\noindent \latin{V. Gloria Patri, et Filio, et Spiritui sancto.}\par

\medskip

\noindent I look from afar, and, behold, I see the Power of God coming, and a cloud covering all the land. Go ye out to meet Him, and say: Tell us if Thou art He, That shalt reign over God’s people Israël.\par
\noindent Verse. Both low and high, rich and poor together.\par
\noindent Answer. Go ye out to meet Him, and say.\par
\noindent Verse. Give ear, O Shepherd of Israël, Thou That leadest Joseph like a flock.\par
\noindent Answer. Tell us if Thou art He.\par
\noindent Verse. Lift up your gates, O ye princes; and be ye lift up, ye everlasting doors, and the King of glory shall come in.\par
\noindent Answer. That shalt reign over God’s people Israël.\par
\noindent Verse. Glory be to the Father, and to the Son, and to the Holy Ghost.\par
\noindent Answer. I look from afar, and, behold, I see the Power of God coming, and a cloud covering all the land. Go ye out to meet Him, and say: Tell us if Thou art He, That shalt reign over God’s people Israël.\par

\subsection{The Entrance Antiphon of the Fourth Sunday}

\latin{Rorate caeli}, from Isaiah 45:8, sung or said as the Mass of the
Fourth Sunday of Advent begins. Its words are those of the Introit of the same
Sunday in the Missal before 1970.

\begin{quote}
\latin{Roráte, cæli, désuper, et nubes pluant justum: aperiátur terra, et gérminet Salvatórem.}\par
\end{quote}

\begin{quote}
Drop down dew, ye heavens, from above, and let the clouds rain the just; let the earth be opened and bud forth a Saviour.\par
\end{quote}

\subsection{\latin{Rorate caeli desuper}}

A chant of Advent, its verses
taken from Isaiah; Dom Prosper Guéranger gives it as the Advent prayer of the
churches of France. It is not a text of the present Missal, though its refrain
is the opening of the Entrance Antiphon above.

\begin{quote}
\latin{Rorate, coeli, desuper, et nubes pluant Justum.}\par
\latin{Ne irascaris, Domine, ne ultra memineris iniquitatis: ecce civitas Sancti facta est deserta, Sion deserta facta est, Jerusalem desolata est, domus sanctificationis nostrae et gloriae tuae, ubi laudaverunt te patres nostri.}\par
\latin{Rorate, coeli, desuper, et nubes pluant Justum.}\par
\latin{Peccavimus, et facti sumus tamquam immundus nos, et cecidimus quasi folium universi; et iniquitates nostrae quasi ventus abstulerunt nos: abscondisti faciem tuam a nobis, et allisisti nos in manu iniquitatis nostrae.}\par
\latin{Rorate, coeli, desuper, et nubes pluant Justum.}\par
\latin{Vide, Domine, afflictionem populi tui, et emitte quem missurus es. Emitte Agnum dominatorem terrae de petra deserti ad montem filiae Sion, ut auferat ipse jugum captivitatis nostrae.}\par
\latin{Rorate, coeli, desuper, et nubes pluant Justum.}\par
\latin{Consolamini, consolamini, popule meus: cito veniet salus tua: quare moerore consumeris? quare comprehendit te dolor? Salvabo te; noli timere: ego enim sum Dominus Deus tuus, Sanctus Israël, Redemptor tuus.}\par
\latin{Rorate, coeli, desuper, et nubes pluant Justum.}\par
\end{quote}

\begin{quote}
Drop down dew, ye heavens, from above, and let the clouds rain the Just One.\par
Be not angry, O Lord, and remember no longer our iniquity: behold the city of thy Sanctuary is become a desert, Sion is made a desert. Jerusalem is desolate, the house of our holiness and of thy glory, where our fathers praised thee.\par
Drop down dew, ye heavens, from above, and let the clouds rain the Just One.\par
We have sinned, and we are become as one unclean, and we have all fallen as a leaf; and our iniquities, like the wind, have taken us away: thou hast hid thy face from us, and hast crushed us by the hand of our iniquity.\par
Drop down dew, ye heavens, from above, and let the clouds rain the Just One.\par
See, O Lord, the affliction of thy people, and send Him whom thou hast promised to send. Send forth the Lamb, the ruler of the earth, from the rock of the desert to the mount of the daughter of Sion, that himself may take off the yoke of our captivity.\par
Drop down dew, ye heavens, from above, and let the clouds rain the Just One.\par
Be comforted, be comforted, my people; thy salvation shall speedily come: why wilt thou waste away in sadness? why hath sorrow seized thee? I will save thee; fear not: for I am the Lord thy God, the Holy One of Israel, thy Redeemer.\par
Drop down dew, ye heavens, from above, and let the clouds rain the Just one.\par
\end{quote}

\subsection{\latin{Veni, Redemptor gentium}}

St Ambrose's hymn for the birth of the Lord, for his day, 7 December, the
day he was ordained bishop. Of the English of the \textit{Hymnal Noted}
(1851), the first, fourth, fifth and seventh stanzas are given.

\begin{verse}
\latin{Veni Redemptor gentium,} \\
\latin{Ostende partum Virginis,} \\
\latin{Miretur omne sæculum,} \\
\latin{Talis decet partus Deum.}

\latin{Non ex virili semine} \\
\latin{Sed mystico spiramine} \\
\latin{Verbum Dei factum est caro,} \\
\latin{Fructusque ventris floruit.}

\latin{Alvus tumescit Virginis,} \\
\latin{Claustra pudoris permanent,} \\
\latin{Vexilla virtutum micant,} \\
\latin{Versatur in templo Deus.}

\latin{Procedens è thalamo suo,} \\
\latin{Pudoris aula regia,} \\
\latin{Geminæ gigas substantiæ} \\
\latin{Alacris ut currat viam.}

\latin{Egressus ejus à Patre,} \\
\latin{Regressus ejus ad Patrem,} \\
\latin{Excursus usque ad inferos,} \\
\latin{Recursus ad sedem Dei.}

\latin{Æqualis æterno Patri,} \\
\latin{Carnis stropheo accingere,} \\
\latin{Infirma nostri corporis} \\
\latin{Virtute firmans perpetim.}

\latin{Præsepe jam fulget tuum,} \\
\latin{Lumenque nox spirat novum,} \\
\latin{Quod nulla nox interpolet,} \\
\latin{Fideque jugi luceat.}

\latin{Deo Patri sit gloria,} \\
\latin{Ejusque soli Filio,} \\
\latin{Cum Spiritu Paraclito,} \\
\latin{Et nunc, et imperpetuum.}
\end{verse}

\begin{verse}
Come, Thou Redeemer of the earth, Come, testify Thy Virgin-birth: \\
All lands admire, all times applaud: Such is the birth that fits a God.

Proceeding from His Chamber free, The royal hall of chastity, \\
Giant of twofold substance, straight His destin'd way He runs elate.

From God the Father He proceeds: To God the Father back He speeds: \\
Proceeds, as far as very hell; Speeds back, to light ineffable.

Thy cradle here shall glitter bright, And darkness breathe a newer light; \\
Where endless faith shall shine serene, And twilight never intervene.
\end{verse}

\subsection{\latin{Alma Redemptoris Mater}}

An antiphon of Our Lady for the end of the day. The Roman Breviary set it
at the close of Compline, the Hour before sleep, from Vespers of the Saturday
before the First Sunday of Advent until the Purification.

\begin{verse}
\latin{Alma Redemptoris Mater, quae pervia coeli} \\
\latin{Porta manes, et stella maris, succurre cadenti,} \\
\latin{Surgere qui curat, populo: tu quae genuisti,} \\
\latin{Natura mirante, tuum sanctum Genitorem,} \\
\latin{Virgo prius ac posterius, Gabrielis ab ore} \\
\latin{Sumens illud Ave, peccatorum miserere.}
\end{verse}

\begin{verse}
Maiden! Mother of Him Who redeemed us, thou that abidest \\
Heaven’s open gate, and the Star of the Sea, come, succour the fallen! \\
Fallen indeed we are, but fain would rise by thy succour. \\
Thou that beyond nature’s course hast borne in time the Eternal; \\
Thou that a Virgin before, and after that childbirth remainest, \\
From the Archangel’s lips the quickening message receiving, \\
Mother of Jesus and us, turn thine eyes of mercy on sinners.
\end{verse}

\subsection{The Canticle of Zachary}

Luke 1:68--79. The Church sings it every morning at Morning Prayer, and on
24 December it is the Gospel of the Mass in the morning.

\begin{quote}
Blessed be the Lord God of Israel: because he hath visited and wrought the redemption of his people.\par
And hath raised up an horn of salvation to us, in the house of David his servant.\par
As he spoke by the mouth of his holy prophets, who are from the beginning.\par
Salvation from our enemies and from the hand of all that hate us.\par
To perform mercy to our fathers and to remember his holy testament.\par
The oath, which he swore to Abraham our father, that he would grant to us.\par
That being delivered from the hand of our enemies, we may serve him without fear:\par
In holiness and justice before him, all our days.\par
And thou, child, shalt be called the prophet of the Highest: for thou shalt, go before the face of the Lord to prepare his ways:\par
To give knowledge of salvation to his people, unto the remission of their sins.\par
Through the bowels of the mercy of our God, in which the Orient from on high hath visited us:\par
To enlighten them that sit in darkness and in the shadow of death: to direct our feet into the way of peace.\par
\end{quote}

\subsection{The Canticle of Mary}

Luke 1:46--55. The Church sings it every evening at Evening Prayer, with the
great antiphons above from 17 to 23 December. On the Third Sunday of Advent in
Year~B it is sung as the responsorial psalm.

\begin{quote}
My soul doth magnify the Lord.\par
And my spirit hath rejoiced in God my Saviour.\par
Because he hath regarded the humility of his handmaid: for behold from henceforth all generations shall call me blessed.\par
Because he that is mighty hath done great things to me: and holy is his name.\par
And his mercy is from generation unto generations, to them that fear him.\par
He hath shewed might in his arm: he hath scattered the proud in the conceit of their heart.\par
He hath put down the mighty from their seat and hath exalted the humble.\par
He hath filled the hungry with good things: and the rich he hath sent empty away.\par
He hath received Israel his servant, being mindful of his mercy.\par
As he spoke to our fathers: to Abraham and to his seed for ever.\par
\end{quote}

\subsection{Before Communion}

In the present Missal priest and people say \latin{Domine, non sum dignus}
together, once, before Communion, and at Mass the words are said in the English
of that Missal. They are the centurion's (Matthew 8:8), who asked Christ to
heal his servant; the prayer asks healing for the soul. The English below is
that of a Missal for the laity printed in 1861.

\begin{quote}
\latin{Domine, non sum dignus ut intres sub tectum meum; sed tantum dic verbo, et sanabitur anima mea.}\par
\end{quote}

\begin{quote}
Lord, I am not worthy that thou shouldst enter under my roof; say but the word, and my soul shall be healed.\par
\end{quote}

\subsection{An Act of Contrition}

For confession before Christmas, and for night prayer after the examination
of conscience.

\begin{quote}
O my God! I am heartily sorry for having offended Thee, and I detest all my sins, because I dread the loss of heaven and the pains of hell, but most of all because they offend Thee, my God, who art all-good and deserving of all my love. I firmly resolve, with the help of Thy grace, to confess my sins, to do penance, and to amend my life. Amen.\par
\end{quote}

\subsection{The St Andrew Prayer}

Said from the feast of St Andrew, 30 November, until Christmas, as the note
printed after it describes.

\begin{quote}
Hail and blessed be the hour and moment in which the Son of God was born of the most pure Virgin Mary at midnight in Bethlehem in the piercing cold. In that hour vouchsafe, O my God! to hear my prayer and grant my desires, through the merits of Our Saviour Jesus Christ, and of His blessed Mother. Amen.\par
It is piously believed that whoever recites the above prayer fifteen times a day, from the feast of St. Andrew (Nov. 30) until Christmas, will obtain what is asked.\par
\end{quote}

\subsection{\latin{O admirabile commercium}}

For the days of Christmas. The Roman Breviary set it as the first
antiphon of Lauds on 1 January, the octave day of Christmas, and took it again
at both Vespers of that day.

\begin{quote}
\latin{O admirabile commercium: * Creator generis humani, animatum corpus sumens, de Virgine nasci dignatus est: et procedens homo sine semine, largitus est nobis suam Deitatem.}\par
\end{quote}

\begin{quote}
How wonderful is the dealing here! * The Maker of man hath taken to Himself a living body, and hath been pleased to be born of the Virgin: He is come forth conceived without seed, and hath made us partakers of His Divine nature.\par
\end{quote}
